\section{Using the audit in an evaluation workflow}
\label{sec:use}
The artifact is intended for researchers and engineering teams evaluating a monitor for an agent with explicit tool-use policies. Its immediate output is an auditable comparison of warnings, not an automatic approval to deploy or block actions. A team can reproduce the saved-record analysis without model access, then adapt the protocol to a separate, locally validated test set.

\paragraph{A concrete decision.}
Suppose a team must decide whether a cancellation agent's monitor produces unnecessary warnings. An executed cancellation, a reference-state match, and a permitted cancellation are different observations. First state which one defines an error. In our final autonomous cohort, an exact interaction-rule detector appears to produce 84/133 false alarms if every state success is declared negative; on the 49 interaction-clean state successes it produces 0/49 by construction. This does not mean that the detector improved. Only the comparison set changed. Nor does conditioning uniformly reduce a model's alert rate: DeepSeek changes from 70/133 to 28/49. A team should examine the selected negative class before replacing a monitor or changing its threshold.

\paragraph{Inputs and sequence.}
For a new application, retain the policy version, tool schemas, ordered public messages and tool responses, proposed actions, task/repeat identifiers, monitor configuration, and independently computed outcome labels. Keep outcome labels and future messages outside monitor inputs. The evaluation can then follow five steps:
\begin{enumerate}
\item \textbf{Declare the property and its observation time.} Separate final-state success from each interaction or eligibility rule. Specify whether the monitor sees a proposed action before execution, or a result afterward. A pre-execution hook is an integration requirement, not an intervention benefit demonstrated by this study.
\item \textbf{Validate the labels and coverage.} Implement deterministic predicates where the rule and facts permit them, check contrasting fixtures, and collect both clean and violating outcomes. If the rule requires contextual judgment, use an adjudication process suited to that rule and report disagreement. Our executable subset does not remove that requirement for broader semantic claims.
\item \textbf{Freeze a matched comparison.} Give each monitor the same permitted prefix and record its exact configuration. Include transparent rule baselines. Separate a change to the public policy, a change to the monitoring question, and a change to the acting agent. Select thresholds using separate development data when calibration is needed.
\item \textbf{Report the denominator and clock.} Report warnings per state success and per measured-rule-clean state success, together with task counts, violation-detection counts and before/at/after timing. A missing output is unknown; a missing negative class is undefined. Preserve interrupted attempts and resample whole tasks rather than treating prefixes as independent observations.
\item \textbf{Choose the next test from the discrepancy.} If a baseline already explains a warning, inspect label alignment before adding model complexity. If a model misses an eligibility violation, inspect its public evidence, target and threshold separately. Evaluate any resulting blocking or repair policy prospectively, measuring task completion, unnecessary intervention, latency and cost on held-out cases.
\end{enumerate}

\paragraph{What a user runs and receives.}
The supplementary \texttt{USAGE.md} provides a short entry point. In a Python environment with the pinned analysis dependencies, \texttt{sh reproduce.sh} regenerates the three cohorts' analyses, manuscript tables and vector figures without model calls. The per-cohort results remain separate. \texttt{replication/analysis.json} contains the expansion's conditional rates, warning timing, repeated-run comparisons and eligibility contrasts; \texttt{replication/evidence\_ledger.json} links reported counts to their evidence. The supplied verifiers check freezes, required scoring cells and request provenance. The usage guide includes a small command that reads these saved results and explains each denominator.

\paragraph{From audit to application.}
This workflow addresses a limited part of risk measurement. The NIST AI RMF calls for selecting appropriate metrics and documenting risks that are not measured \citep{nist2023rmf}; its generative-AI profile cautions against extrapolating capabilities from narrow assessments \citep{nist2024genai}. Our audit is compatible with those reporting principles, but is not a NIST certification or a complete risk-management process. Teams would still need application-specific policy coverage, representative tests, security assessment and prospective intervention evidence before relying on a monitor operationally.
